\chapter*{Prefaci}
\addcontentsline{toc}{chapter}{Prefaci}

Aquest llibre és una introducció a la \textbf{teoria de categories} pensada per a un primer contacte, sense pressuposar cap coneixement previ de la matèria. El llenguatge de les categories unifica construccions que apareixen una vegada i una altra en àrees molt diverses de les matemàtiques, de la informàtica i de moltes altres ciències, i ensenya a estudiar les estructures no només per la seva constitució interna, sinó també per la manera com es relacionen entre elles.

El llibre segueix l'estructura habitual dels materials d'\textbf{aprendes}. Després d'una introducció general, el contingut es distribueix en quatre parts complementàries. La \textbf{Teoria} presenta les definicions, els resultats i els exemples representatius; la \textbf{Pràctica} reprèn els mateixos capítols mitjançant exercicis resolts amb detall; els \textbf{Exercicis proposats} permeten treballar autònomament cada tema i inclouen indicacions breus; finalment, els \textbf{Tests interactius} serveixen per comprovar la comprensió dels conceptes essencials.

La intenció és que aquestes quatre parts es puguin utilitzar de manera successiva: primer entendre els conceptes, després observar com s'usen en problemes concrets, posar-los en pràctica sense una solució completa i, finalment, comprovar-ne l'assimilació. Els exemples de la part teòrica no substitueixen la pràctica: tenen la funció de donar significat a les definicions i mostrar-ne les primeres manifestacions.

El recorregut va de les nocions bàsiques ---categories, functors i transformacions naturals--- fins als classificadors de subobjectes i una primera mirada als topos, passant per Yoneda, límits, adjuncions i categories cartesianes tancades. La lògica categòrica pròpiament dita ---sistemes deductius, doctrines i semàntica--- es reserva per a un volum posterior, del qual aquest text vol ser el fonament.

Al llarg del llibre donarem un paper especial als diagrames: quan una propietat categòrica tingui una expressió natural mitjançant fletxes, procurarem mostrar-la juntament amb la seva formulació algebraica. També tornarem sistemàticament, sempre que sigui pertinent, a alguns exemples de referència ---conjunts, preordres, lògica i grafs--- per aprendre a reconèixer una mateixa estructura sota formes matemàtiques diferents. 